%% kapa: exam=1 round=33 year=2022 period=2 form=A subject=accounting subject_ko=회계학 seq=5 number=45
%% source: Q-Net 감정평가사 기출문제 (cst00309) artlSeq=5209997 file=raw/1008/5209997/2022년도 제33회 감정평가사 1차시험 2교시(감정평가 관계법규, 회계학).pdf
%% conversion: kordoc 4.24.51 -> tools/gen1.py (auto); review=proofread-sonnet; flags=none
\begin{question}{45}
재무제표 요소의 측정기준에 관한 설명으로 옳은 것은?
\begin{choices}
  \item 공정가치는 측정일 현재 동등한 자산의 원가로서 측정일에 지급할 대가와 그 날에 발생할 거래원가를 포함한다.
  \item 현행원가는 자산을 취득 또는 창출할 때 발생한 원가의 가치로서 자산을 취득 또는 창출하기 위하여 지급한 대가와 거래원가를 포함한다.
  \item 사용가치는 기업이 자산의 사용과 궁극적인 처분으로 얻을 것으로 기대하는 현금흐름 또는 그 밖의 경제적효익의 현재가치이다.
  \item 이행가치는 측정일에 시장참여자 사이의 정상거래에서 부채를 이전할 때 지급하게 될 가격이다.
  \item 역사적 원가는 측정일 현재 자산의 취득 또는 창출을 위해 이전해야 하는 현금이나 그 밖의 경제적 자원의 현재가치이다.
\end{choices}
\end{question}
\answer{3}
